\scanpage{215}

\mspage{89}

Sous ces conditions, \unsure{\ldots|on a}

\[
\underbrace{\mathrm{Mon}}_{C}\cap W=\widetilde{\mathrm{Mon}\cap W}\qquad
\begin{aligned}&\text{i.e. }\mathrm{Mon}\cap W\text{ est}\\&\quad Q\text{-clos à gauche},\end{aligned}
\tag*{(7.2.\deleted{3}4)}
\]

et

\[
\begin{aligned}
&\text{Il existe une petite partie }\mathrm{TC}_0\subset\mathrm{Mon}\cap W\\
&\text{telle que }\mathrm{Mon}\cap W=\widetilde{\mathrm{TC}_0}.
\end{aligned}\tag*{(7.2.\deleted{\unsure{4}}5)}
\]

[De façon plus précise, si $\Pi\geq\unsure{\Pi_0|k}$ est un cardinal tel que $W$ soit $\Pi$-accessible, on peut prendre \deleted{pour} $\mathrm{TC}_0=(\mathrm{Mon}\cap W)^{\Pi}$ (ou une partie \uline{petite} de $(\mathrm{Mon}\cap W)^{\unsure{\Pi|P}}$ ayant les \unsure{$=^{\mathrm{s}}$} classes d'isomorphie).] \deleted{\unsure{pour t}} Donc (th. 1 du \S\,1) le couple $(\mathrm{TC}=\mathrm{Mon}\cap W,\ F=(\mathrm{TC})_*)$ est un couple de Quillen clos, et il y a $=$ une factorisation \uline{fonctorielle} de toute flèche $f$ en

\[
f=pi\ ;\quad i\in\unsure{\mathrm{TC}|\mathrm{TNC}},\ p\unsure{\in|\vdash}F.\tag*{(7.2.\unsure{4'|4}.)}
\]

\deleted{Mieux}, Si on pose

\[
C=\mathrm{Mon},\ \mathrm{TC}=\mathrm{Mon}\cap W,\ F=(\mathrm{TC})_*,\ \mathrm{TF}=C_*,\tag*{(7.2.5.)}
\]

le triple $(W,C,F)$ est un triple de Quillen clos, sss on \deleted{\unsure{suppose que l'on a}} \inserted{a les conditions \inserted{suppl.} suivantes}

$\left\{\vphantom{\dfrac{A}{A}}\right.$ (7.2.6) \textcircled{g} \deleted{\illegible{1}} (sur $M$ seulemt) $\mathrm{Mon}$ stable par cochg. de base.

\[
\underset{\substack{\Vert\,\mathrm{def}\\(\mathrm{Mon})_*}}{\mathrm{TF}}\subset W\deleted{\unsure{\mathrm{univ}}} \tag*{(7.2.\deleted{\unsure{6}}7)}
\]

i.e. si \inserted{d)} \unsure{toute} $f\colon X\mathrel{-}Y$ qui fait de $X$ un objet \uline{injectif} de $M_{/Y}$, est dans $W$ (donc aussi dans $W_{\mathrm{univ}}$).
% archive sheet number: 214
